% Curso intensivo del Capítulo 12 de la tesis (corrección del algoritmo
% Quicksort) y del Apéndice B.2 (conservación mediante conteos).
% Escrito en matemática de papel, sin Lean. Fuente:
%   tesis-mate/documento/capitulos/5_practica.tex
%   tesis-mate/documento/capitulos/apendice_aristotle.tex
% Compilar dos veces:  lualatex CursoQuicksort.tex
\documentclass[11pt,a4paper]{article}
\usepackage[a4paper,margin=2.3cm,headheight=14pt]{geometry}
\usepackage{amsmath,mathtools}
\usepackage{fontspec}
\usepackage{unicode-math}
\usepackage[spanish,es-noshorthands,es-tabla]{babel}
\setmainfont{TeX Gyre Pagella}
\setmathfont{TeX Gyre Pagella Math}
\setsansfont{DejaVu Sans}[Scale=MatchLowercase]
\usepackage{xcolor}
\usepackage[most]{tcolorbox}
\usepackage{booktabs,tabularx,array,longtable}
\usepackage{tikz}
\usetikzlibrary{arrows.meta,positioning,trees,fit}
\usepackage{enumitem}
\usepackage{fancyhdr}
\usepackage[hidelinks]{hyperref}
\usepackage{microtype}

\definecolor{fondores}{HTML}{F4F7FB}
\definecolor{borderes}{HTML}{3F6FA8}
\definecolor{fondonota}{HTML}{F1EEF8}
\definecolor{bordenota}{HTML}{8C7DB8}
\definecolor{fondoejemplo}{HTML}{F7F8FA}
\definecolor{bordeejemplo}{HTML}{B8C2CF}
\definecolor{fondotrampa}{HTML}{FFF4EC}
\definecolor{bordetrampa}{HTML}{D98A4E}
\definecolor{fondoclave}{HTML}{FFF9E6}
\definecolor{bordeclave}{HTML}{C9A227}
\definecolor{fondoesquema}{HTML}{EEF5EE}
\definecolor{bordeesquema}{HTML}{5E9A62}

\newtcolorbox{resultado}[1]{enhanced, breakable, colback=fondores,
  colframe=borderes, boxrule=0.6pt, arc=3pt, colbacktitle=borderes,
  coltitle=white, fonttitle=\sffamily\bfseries, title={#1},
  before skip=10pt, after skip=10pt}
\newtcolorbox{esquema}[1]{enhanced, breakable, colback=fondoesquema,
  colframe=bordeesquema, boxrule=0.6pt, arc=3pt,
  colbacktitle=bordeesquema, coltitle=white, fonttitle=\sffamily\bfseries,
  title={#1}, before skip=10pt, after skip=10pt}
\newtcolorbox{ejemplo}[1]{enhanced, breakable, colback=fondoejemplo,
  colframe=bordeejemplo, boxrule=0.5pt, arc=2pt, colbacktitle=bordeejemplo!40,
  coltitle=black, fonttitle=\sffamily\small\bfseries, title={#1},
  before skip=6pt, after skip=8pt}
\newtcolorbox{nota}[1][Nota]{enhanced, breakable, colback=fondonota,
  colframe=bordenota, boxrule=0.5pt, arc=3pt, colbacktitle=bordenota,
  coltitle=white, fonttitle=\sffamily\bfseries\small, title={#1},
  before skip=8pt, after skip=8pt}
\newtcolorbox{trampa}[1][Lo que pueden preguntar]{enhanced, breakable,
  colback=fondotrampa, colframe=bordetrampa, boxrule=0.5pt, arc=3pt,
  colbacktitle=bordetrampa, coltitle=white, fonttitle=\sffamily\bfseries\small,
  title={#1}, before skip=8pt, after skip=8pt}
\newtcolorbox{clave}{enhanced, colback=fondoclave, colframe=bordeclave,
  boxrule=0.8pt, arc=3pt, left=8pt, right=8pt, top=6pt, bottom=6pt,
  before skip=10pt, after skip=10pt}

\newcommand{\campo}[1]{\par\smallskip\noindent\textbf{#1}\enspace}
\newcommand{\enunciado}{\campo{Enunciado.}}
\newcommand{\palabras}{\campo{En palabras.}}
\newcommand{\sirve}{\campo{Para qué sirve.}}
\newcommand{\idea}{\campo{Idea.}}
\newcommand{\prueba}{\campo{Demostración.}}

% Notación de la tesis (Capítulo 12)
\newcommand{\Perm}{\sim}
\newcommand{\concat}{\mathbin{+\!\!+}}
\newcommand{\qs}{\operatorname{\mathsf{qsort}}}
\newcommand{\filt}{\operatorname{\mathsf{filtrar}}}
\newcommand{\len}[1]{\lvert #1\rvert}
\newcommand{\Sorted}{\operatorname{\mathsf{Ordenada}}}
\newcommand{\Pairwise}{\operatorname{\mathsf{PorPares}}}
\newcommand{\cuenta}{\operatorname{\mathsf{cuenta}}}
\newcommand{\decidir}{\operatorname{\mathsf{decidir}}}
\newcommand{\V}{\mathsf{V}}
\newcommand{\F}{\mathsf{F}}

\pagestyle{fancy}
\fancyhf{}
\fancyhead[L]{\small\sffamily Capítulo 12: corrección de Quicksort}
\fancyhead[R]{\small\sffamily\nouppercase{\leftmark}}
\fancyfoot[C]{\small\thepage}
\renewcommand{\headrulewidth}{0.4pt}

\setlist{itemsep=2pt, topsep=4pt}
\setlength{\parskip}{4pt}
\setlength{\parindent}{0pt}

\title{\textbf{Curso intensivo del Capítulo 12}\\[4pt]
  \Large Corrección del algoritmo Quicksort\\[2pt]
  \large Qué se demuestra, cómo se demuestra y cómo explicarlo}
\author{Material de apoyo para la defensa}
\date{}

\begin{document}
\maketitle
\thispagestyle{empty}

\begin{abstract}
\noindent
Este documento explica el Capítulo 12 de la tesis y la variante del
Apéndice B.2, en matemática de papel. No usa código. Cada resultado aparece en
una ficha con la misma estructura: enunciado, qué dice en palabras, para qué
sirve, la idea en una frase y la demostración paso a paso. Incluye ejemplos
trazados a mano, qué propiedades del orden se usan y dónde, y las preguntas
más probables del tribunal con sus respuestas. La numeración (Lema 12.5,
Teorema 12.11, etc.) es la de la versión compilada de la tesis.
\end{abstract}

\tableofcontents
\newpage

%=====================================================================
\section{El capítulo en una página}
%=====================================================================

\begin{clave}
\centering
\textbf{Teorema principal (Teorema 12.12).}\quad
Para toda lista \(\ell\) sobre un tipo linealmente ordenado,
\[
\qs(\ell)\Perm\ell\qquad\text{y}\qquad\Sorted(\qs(\ell)).
\]
Quicksort devuelve una \emph{permutación ordenada} de su entrada.
\end{clave}

\textbf{El algoritmo.} Para ordenar \(p::r\) (pivote \(p\), cola \(r\)) se
separa la cola en \(L\), los elementos \(\leq p\), y \(R\), los elementos
\(>p\). Se ordenan \(L\) y \(R\) recursivamente y se pegan con el pivote en
medio:
\[
\qs([\,])=[\,],\qquad
\qs(p::r)=\qs(L)\concat\big(p::\qs(R)\big).
\]
Ejemplo: \([3,2,5,1,4]\ \to\ ([2,1],\,3,\,[5,4])\ \to\ [1,2]\concat(3::[4,5])=[1,2,3,4,5]\).

\textbf{Qué hay que demostrar: tres cosas.}
\begin{center}
\renewcommand{\arraystretch}{1.3}
\begin{tabular}{@{}lll@{}}
\toprule
\textbf{Obligación} & \textbf{Por qué hace falta} & \textbf{Resultado}\\
\midrule
\emph{Terminación} & sin ella \(\qs\) no es una función & Proposición 12.6\\
\emph{Contenido}: \(\qs(\ell)\Perm\ell\) & no se pierden ni se inventan elementos & Teorema 12.8\\
\emph{Orden}: \(\Sorted(\qs(\ell))\) & la salida queda en orden no decreciente & Teorema 12.11\\
\bottomrule
\end{tabular}
\end{center}
Las tres se demuestran con la misma estrategia: inducción sobre la
\emph{longitud} de la lista, porque las llamadas recursivas reciben \(L\) y
\(R\), que son más cortas que \(p::r\) pero no son su cola.

\subsection*{Mapa de dependencias}

\begin{center}
\begin{tikzpicture}[font=\footnotesize,
  res/.style={draw, rounded corners, fill=blue!6, align=center, inner sep=3pt, text width=27mm},
  dep/.style={-{Stealth}, thick, gray!60!black}]
\node[res] (fle)   at (0,0)     {Lema 12.5\\filtrar no alarga};
\node[res] (fperm) at (4.2,0)   {Lema 12.7\\partición \(\Perm\) lista};
\node[res] (pw)    at (8.6,0)   {Lema 12.10\\\(\Pairwise\) sobre \(\concat\)};
\node[res] (term)  at (0,-1.8)  {Proposición 12.6\\terminación};
\node[res] (perm)  at (4.2,-1.8){Teorema 12.8\\permutación};
\node[res] (mem)   at (4.2,-3.6){Corolario 12.9\\pertenencia};
\node[res] (sorted)at (8.6,-3.6){Teorema 12.11\\ordenamiento};
\node[res, fill=yellow!20] (corr) at (6.4,-5.4) {\textbf{Teorema 12.12}\\corrección};
\draw[dep] (fle) -- (term);
\draw[dep] (fle) -- (perm);
\draw[dep] (fperm) -- (perm);
\draw[dep] (perm) -- (mem);
\draw[dep] (mem) -- (sorted);
\draw[dep] (pw) -- (sorted);
\draw[dep] (fle.south) .. controls (1.2,-4.2) and (5.5,-3.9) .. (sorted.west);
\draw[dep] (perm.south west) .. controls (2.6,-3.6) and (3.2,-5.4) .. (corr.west);
\draw[dep] (sorted) -- (corr);
\end{tikzpicture}
\end{center}

\subsection*{Todos los resultados de un vistazo}

{\small
\begin{longtable}{@{}p{0.15\textwidth}p{0.47\textwidth}p{0.32\textwidth}@{}}
\toprule
\textbf{Resultado} & \textbf{Enunciado} & \textbf{Técnica}\\
\midrule
\endhead
Lema 12.5 & \(\len{\filt(q)(\ell)}\leq\len\ell\) & inducción sobre la lista\\
Proposición 12.6 & \(\qs\) termina (es total) & inducción bien fundada sobre \(\len\ell\)\\
Lema 12.7 & \(\filt(q)(\ell)\concat\filt(q^c)(\ell)\Perm\ell\) & inducción sobre la lista\\
Teorema 12.8 & \(\qs(\ell)\Perm\ell\) & inducción bien fundada + cadena de permutaciones\\
Corolario 12.9 & \(a\in\qs(\ell)\iff a\in\ell\) & consecuencia de 12.8\\
Lema 12.10 & \(\Pairwise\) de una concatenación & inducción sobre la primera lista\\
Teorema 12.11 & \(\Sorted(\qs(\ell))\) & inducción bien fundada + 12.10\\
Teorema 12.12 & permutación y orden & conjunción de 12.8 y 12.11\\
Apéndice B.2 & otra prueba de 12.8 contando apariciones & inducción + aritmética\\
\bottomrule
\end{longtable}
}

%=====================================================================
\section{El escenario: definiciones}
%=====================================================================

\subsection{El orden}

\(\alpha\) es un tipo con un orden lineal \(\leq\): reflexivo
(\(a\leq a\)), antisimétrico, transitivo y total
(\(a\leq b\) o \(b\leq a\)). El orden estricto cumple
\[
a<b\iff(a\leq b\wedge a\neq b),\qquad a<b\Rightarrow a\leq b,\qquad
\neg(b\leq a)\iff a<b .\tag{O}
\]
La última ley, consecuencia del orden lineal, dice que lo que \emph{no}
es \(\leq p\) es \(>p\). Es lo que hace que la partición no deje nada fuera.

\subsection{Listas y operaciones}

\begin{itemize}
  \item \(a::r\) es la lista con cabeza \(a\) y cola \(r\); \([\,]\) es la
    vacía; \(\concat\) es la concatenación; \(\len\ell\) es la longitud.
  \item Pertenencia: \(x\in a::r\iff x=a\vee x\in r\), y
    \(x\in\ell_1\concat\ell_2\iff x\in\ell_1\vee x\in\ell_2\).
  \item \(\cuenta(a,\ell)\): número de apariciones de \(a\) en \(\ell\).
  \item \textbf{Filtro (F).} \(\filt(q)\) conserva, en su orden, los elementos
    con \(q(x)=\V\):
    \[
    \filt(q)([\,])=[\,],\qquad
    \filt(q)(a::r)=\begin{cases}a::\filt(q)(r)&\text{si }q(a)=\V,\\ \filt(q)(r)&\text{si }q(a)=\F.\end{cases}
    \]
    En consecuencia, \(x\in\filt(q)(\ell)\iff x\in\ell\wedge q(x)=\V\).
  \item \(q^c\) es el complemento booleano de \(q\): vale \(\V\) donde \(q\)
    vale \(\F\), y viceversa.
\end{itemize}

\subsection{Las dos mitades de la especificación}

\begin{itemize}
  \item \textbf{Ordenada (Def.~12.2).} \(\Sorted(\ell):=\Pairwise(\leq)(\ell)\),
    donde
    \[
    \Pairwise\,R\,[\,]\ \text{siempre},\qquad
    \Pairwise\,R\,(a::r)\iff(\forall x\in r,\ R\,a\,x)\wedge\Pairwise\,R\,r .\tag{P}
    \]
    En palabras: cada elemento está relacionado con \emph{todos} los que van
    detrás de él. Equivale a la caracterización por índices
    \(\forall\,i<j,\ a_i\leq a_j\).
  \item \textbf{Permutación (Def.~12.3).} \(\ell_1\Perm\ell_2\) significa
    que difieren solo en el orden:
    \(\forall a,\ \cuenta(a,\ell_1)=\cuenta(a,\ell_2)\). Se usan estas
    reglas:
    \begin{center}
    {\small
    \begin{tabular}{@{}ll@{}}
    \toprule
    reflexiva, simétrica, transitiva & \(\Perm\) es de equivalencia\\
    congruencia de \(\concat\) & \(\ell_1\Perm\ell_1'\wedge\ell_2\Perm\ell_2'\Rightarrow\ell_1\concat\ell_2\Perm\ell_1'\concat\ell_2'\)\\
    congruencia de \(::\) & \(\ell_1\Perm\ell_2\Rightarrow a::\ell_1\Perm a::\ell_2\)\\
    elemento intermedio & \(L\concat(a::R)\Perm a::(L\concat R)\)\\
    invariancia de pertenencia & \(\ell_1\Perm\ell_2\Rightarrow(a\in\ell_1\iff a\in\ell_2)\)\\
    \bottomrule
    \end{tabular}}
    \end{center}
\end{itemize}

\begin{ejemplo}{Por qué la especificación pide las dos cosas (§12.2.1)}
\begin{itemize}
  \item \(f_0(\ell):=[\,]\) siempre da una lista ordenada, pero
    \(f_0([1])=[\,]\not\Perm[1]\): pierde elementos.
  \item \(f_{\mathrm{id}}(\ell):=\ell\) siempre da una permutación, pero
    \(f_{\mathrm{id}}([2,1])=[2,1]\) no está ordenada.
\end{itemize}
Ninguna de las dos condiciones basta sola.
\end{ejemplo}

\begin{ejemplo}{Por qué la permutación se mide con multiplicidades}
\([1,1]\) y \([1]\) tienen los mismos elementos, pero no son permutación una
de otra: \(\cuenta(1,[1,1])=2\neq1\). Por eso el Corolario 12.9 (misma
pertenencia) es \emph{más débil} que el Teorema 12.8 (permutación). Una
función que borrara duplicados cumpliría 12.9 pero no 12.8.
\end{ejemplo}

\subsection{El algoritmo (Def.~12.4)}

Con los predicados \(q_{\leq p}(x):=\decidir(x\leq p)\) y
\(q_{>p}:=q_{\leq p}^c\),
\[
L:=\filt(q_{\leq p})(r),\qquad R:=\filt(q_{>p})(r),\qquad
\qs(p::r)=\qs(L)\concat(p::\qs(R)).
\]
Por el filtro y la ley (O):
\[
x\in L\iff x\in r\wedge x\leq p,\qquad
x\in R\iff x\in r\wedge p<x .
\]

%=====================================================================
\section{El algoritmo en acción}
%=====================================================================

\begin{ejemplo}{Traza de \(\qs([3,1,2,5,4,9])\)}
\begin{center}
\begin{tikzpicture}[font=\small,level distance=12mm,
  level 1/.style={sibling distance=46mm},
  level 2/.style={sibling distance=23mm},
  every node/.style={draw,circle,inner sep=1.4pt,minimum size=7mm,fill=blue!8},
  tam/.style={draw=none,fill=none,font=\scriptsize,text=black!70},
  et/.style={draw=none,fill=none}]
  \node[label={[tam]right:$n=6$}] {3}
    child { node[label={[tam]left:$n=2$}] {1}
      child { node[fill=green!12] {} edge from parent node[et,left]{$\leq$} }
      child { node[label={[tam]right:$n=1$}] {2}
        child { node[fill=green!12] {} edge from parent node[et,left]{$\leq$}}
        child { node[fill=green!12] {} edge from parent node[et,right]{$>$}}
        edge from parent node[et,right]{$>$} }
    }
    child { node[label={[tam]right:$n=3$}] {5}
      child { node[label={[tam]right:$n=1$}] {4}
        child { node[fill=green!12] {} edge from parent node[et,left]{$\leq$}}
        child { node[fill=green!12] {} edge from parent node[et,right]{$>$}}
        edge from parent node[et,left]{$\leq$} }
      child { node[label={[tam]right:$n=1$}] {9}
        child { node[fill=green!12] {} edge from parent node[et,left]{$\leq$}}
        child { node[fill=green!12] {} edge from parent node[et,right]{$>$}}
        edge from parent node[et,right]{$>$} }
    };
\end{tikzpicture}
\end{center}
\begin{center}
{\footnotesize
\renewcommand{\arraystretch}{1.25}\setlength{\tabcolsep}{4pt}
\begin{tabular}{@{}lllll@{}}
\toprule
\textbf{Llamada} & \textbf{Pivote} & \(L\) (\(\leq p\)) & \(R\) (\(>p\)) & \textbf{Resultado}\\
\midrule
\(\qs([3,1,2,5,4,9])\) & 3 & \([1,2]\) & \([5,4,9]\) & \([1,2]\concat(3::[4,5,9])\)\\
\(\qs([1,2])\) & 1 & \([\,]\) & \([2]\) & \([\,]\concat(1::[2])=[1,2]\)\\
\(\qs([5,4,9])\) & 5 & \([4]\) & \([9]\) & \([4]\concat(5::[9])=[4,5,9]\)\\
\(\qs([2])\), \(\qs([4])\), \(\qs([9])\) & --- & \([\,]\) & \([\,]\) & la propia lista\\
\bottomrule
\end{tabular}}
\end{center}
La primera fila da \([1,2]\concat(3::[4,5,9])=[1,2,3,4,5,9]\).
Las longitudes \(n\) muestran la terminación: cada hijo es estrictamente más
corto que su padre (\(6\to2,3\); \(2\to0,1\); \(3\to1,1\)). Leer el árbol de
izquierda a derecha (recorrido simétrico) da la salida ordenada.
\end{ejemplo}

\begin{ejemplo}{Elementos repetidos: \(\qs([2,1,2,1])\)}
Pivote \(2\), cola \([1,2,1]\): \(L=[1,2,1]\) (todos \(\leq2\)) y \(R=[\,]\).
\[
\qs([1,2,1]):\ \text{pivote }1,\ L=[1],\ R=[2]\ \Rightarrow\ [1]\concat(1::[2])=[1,1,2].
\]
\[
\qs([2,1,2,1])=[1,1,2]\concat(2::[\,])=[1,1,2,2].
\]
Las copias del pivote van a la izquierda, porque el criterio es \(\leq\). Las
multiplicidades se conservan: \(\cuenta(1)=2\) y \(\cuenta(2)=2\) antes y
después.
\end{ejemplo}

%=====================================================================
\section{Las tres técnicas del capítulo}
%=====================================================================

\begin{esquema}{Técnica 1 · Inducción estructural sobre listas}
Para probar \(\forall\ell.\ P(\ell)\): se prueba \(P([\,])\) y, suponiendo
\(P(r)\) (hipótesis inductiva sobre la cola), se prueba \(P(a::r)\). Casi
siempre el caso \(a::r\) se divide según el valor de \(q(a)\), porque así
está definido el filtro.

\textbf{Dónde:} Lemas 12.5, 12.7 y 12.10.
\end{esquema}

\begin{esquema}{Técnica 2 · Inducción bien fundada sobre la longitud}
Para probar \(\forall\ell.\ P(\ell)\) se puede suponer \(P(s)\) para
\emph{toda} lista \(s\) con \(\len s<\len\ell\), y no solo para la cola.

\textbf{Por qué hace falta.} Las llamadas recursivas de \(\qs(p::r)\) son
sobre \(L\) y \(R\), no sobre \(r\). La inducción estructural solo daría la
hipótesis para \(r\). Como \(\len L,\len R\leq\len r<\len{p::r}\), la
inducción por longitud da la hipótesis justo donde hace falta.

\textbf{Por qué es válida.} \((\mathbb N,<)\) no tiene cadenas descendentes
infinitas; es la inducción fuerte sobre los naturales.

\textbf{Plantilla.} «Caso \([\,]\): directo. Caso \(p::r\): por el
Lema 12.5, \(\len L<\len{p::r}\) y \(\len R<\len{p::r}\); por hipótesis
inductiva, \(P(L)\) y \(P(R)\); entonces \dots»

\textbf{Dónde:} Proposición 12.6 y Teoremas 12.8 y 12.11. Las tres pruebas
empiezan igual.
\end{esquema}

\begin{esquema}{Técnica 3 · Cadenas de permutaciones y de equivalencias}
Una permutación entre dos listas lejanas se prueba con una cadena
\(\ell_0=\ell_1\Perm\ell_2\Perm\cdots\Perm\ell_n\), con una regla de la
Definición 12.3 en cada eslabón; la transitividad cierra la cadena. Del mismo
modo, la equivalencia del Lema 12.10 se prueba con una cadena de \(\iff\).

\textbf{Dónde:} Lema 12.7, Teorema 12.8 y Lema 12.10.
\end{esquema}

%=====================================================================
\section{Los resultados, uno por uno}
%=====================================================================

\subsection{Terminación}

%---------------------------------------------------------------------
\begin{resultado}{Lema 12.5 · Filtrar no alarga la lista}
\enunciado Para todo predicado \(q\) y toda lista \(\ell\),
\(\len{\filt(q)(\ell)}\leq\len\ell\).

\palabras Quitar elementos no puede hacer una lista más larga.

\sirve Da las cotas \(\len L,\len R\leq\len r\) que se usan en las tres
inducciones por longitud.

\prueba Por inducción sobre \(\ell\).
\begin{itemize}
  \item \(\ell=[\,]\): \(0\leq0\).
  \item \(\ell=a::r\), con hipótesis \(\len{\filt(q)(r)}\leq\len r\):
    \begin{itemize}
      \item si \(q(a)=\V\):
        \(\len{\filt(q)(a::r)}=1+\len{\filt(q)(r)}\leq1+\len r=\len{a::r}\);
      \item si \(q(a)=\F\):
        \(\len{\filt(q)(a::r)}=\len{\filt(q)(r)}\leq\len r<\len{a::r}\).
    \end{itemize}
\end{itemize}
\end{resultado}

%---------------------------------------------------------------------
\begin{resultado}{Proposición 12.6 · Terminación}
\enunciado Cada llamada recursiva de \(\qs(p::r)\) recibe una lista
estrictamente más corta; por tanto \(\qs\) es una función total.

\idea La longitud baja en cada llamada, y una sucesión de naturales no puede
bajar para siempre.

\prueba Por inducción bien fundada sobre \(\len\ell\).
\begin{itemize}
  \item \(\ell=[\,]\): \(\qs([\,])=[\,]\), sin llamadas.
  \item \(\ell=p::r\): por el Lema 12.5,
    \(\len L\leq\len r<\len r+1=\len{p::r}\), y lo mismo para \(R\). Por
    hipótesis inductiva \(\qs(L)\) y \(\qs(R)\) están definidas; \(\filt\) y
    \(\concat\) son totales, así que la ecuación define \(\qs(p::r)\).
\end{itemize}
\end{resultado}

\begin{trampa}
\begin{itemize}
  \item \emph{¿Por qué hay que probar la terminación?} En un sistema lógico,
    una «definición» por ecuaciones que no termine puede ser contradictoria.
    Por ejemplo, \(g(n)=g(n)+1\) obligaría a \(0=1\). Probar que la medida
    \(\len\ell\) decrece es lo que garantiza que existe una única función
    total con esas ecuaciones.
  \item \emph{¿Por qué no basta la recursión estructural?} Porque \(L\) y
    \(R\) no son la cola \(r\), sino listas filtradas a partir de ella.
  \item \emph{¿Qué medida se usa?} \(\mu(\ell)=\len\ell\), con valores en
    \((\mathbb N,<)\), que es bien fundado. Thompson usa un argumento
    equivalente con un parámetro \(n\geq\len\ell\) explícito.
\end{itemize}
\end{trampa}

\subsection{Corrección, parte I: la salida es una permutación}

%---------------------------------------------------------------------
\begin{resultado}{Lema 12.7 · La partición por filtro es una permutación}
\enunciado Para todo predicado \(q\) y toda lista \(\ell\),
\[
\filt(q)(\ell)\concat\filt(q^c)(\ell)\ \Perm\ \ell .
\]

\palabras Separar una lista en «los que cumplen \(q\)» y «los que no» y pegar
las dos partes da los mismos elementos, quizá en otro orden.

\sirve Es el hecho combinatorio del paso de partición (con \(q=q_{\leq p}\)):
\(L\concat R\Perm r\).

\idea Cada elemento va exactamente a una de las dos partes, porque \(q\) y
\(q^c\) nunca valen \(\V\) a la vez.

\prueba Por inducción sobre \(\ell\). Sea \(L:=\filt(q)(r)\) y
\(R:=\filt(q^c)(r)\); la hipótesis inductiva es \(L\concat R\Perm r\), y por
congruencia de \(::\), \(a::(L\concat R)\Perm a::r\).
\begin{itemize}
  \item \(\ell=[\,]\): \([\,]\concat[\,]=[\,]\Perm[\,]\).
  \item \(\ell=a::r\) con \(q(a)=\V\) (luego \(q^c(a)=\F\)):
    \[
    (a::L)\concat R=a::(L\concat R)\Perm a::r .
    \]
  \item \(\ell=a::r\) con \(q(a)=\F\) (luego \(q^c(a)=\V\)):
    \[
    L\concat(a::R)\ \Perm\ a::(L\concat R)\ \Perm\ a::r ,
    \]
    por la regla del elemento intermedio y la congruencia de \(::\).
\end{itemize}
\end{resultado}

%---------------------------------------------------------------------
\begin{resultado}{Teorema 12.8 · La salida es una permutación de la entrada}
\enunciado Para toda lista \(\ell\), \(\qs(\ell)\Perm\ell\).

\idea Por inducción, cada mitad ordenada es permutación de su mitad; la
partición es permutación de la cola; se recolocan las piezas.

\prueba Por inducción bien fundada sobre \(\len\ell\).
\begin{itemize}
  \item \(\ell=[\,]\): \(\qs([\,])=[\,]\Perm[\,]\).
  \item \(\ell=p::r\): por la Técnica 2, la hipótesis inductiva da
    \(\qs(L)\Perm L\) y \(\qs(R)\Perm R\). El Lema 12.7 (con
    \(q=q_{\leq p}\), cuyo complemento es \(q_{>p}\)) da \(L\concat R\Perm r\).
    Entonces
    \[
    \begin{aligned}
    \qs(p::r)&=\qs(L)\concat(p::\qs(R)) &&\text{definición}\\
    &\Perm L\concat(p::R) &&\text{congruencias de }\concat\text{ y }::\\
    &\Perm p::(L\concat R) &&\text{elemento intermedio}\\
    &\Perm p::r &&\text{congruencia de }::\text{ y Lema 12.7.}
    \end{aligned}
    \]
\end{itemize}
\end{resultado}

%---------------------------------------------------------------------
\begin{resultado}{Corolario 12.9 · Se preserva la pertenencia}
\enunciado \(a\in\qs(\ell)\iff a\in\ell\).

\prueba El Teorema 12.8 y la invariancia de pertenencia de \(\Perm\).

\sirve En el Teorema 12.11, para saber qué elementos hay en \(\qs(L)\) y en
\(\qs(R)\), y así obtener sus cotas respecto del pivote.
\end{resultado}

\begin{nota}[Una prueba alternativa: Apéndice B.2]
En lugar de manipular permutaciones, se puede contar. Con
\(c_x(\ell):=\cuenta(x,\ell)\) y \(\delta_{x,p}=1\) si \(x=p\) y \(0\) si no:
\[
\begin{aligned}
c_x(\qs(p::r))&=c_x(\qs(L))+\delta_{x,p}+c_x(\qs(R)) &&\text{(a)}\\
&=c_x(L)+\delta_{x,p}+c_x(R) &&\text{(b)}\\
&=\delta_{x,p}+c_x(r)=c_x(p::r) &&\text{(c)}
\end{aligned}
\]
(a) el conteo es aditivo sobre \(\concat\) y \(::\); (b) hipótesis
inductiva; (c) cada aparición de la cola está en una sola mitad.
La igualdad \(c_x(L)+c_x(R)=c_x(r)\) es la versión numérica del Lema 12.7.
Ambas pruebas son la misma idea: la partición no pierde ni duplica.
\end{nota}

\subsection{Corrección, parte II: la salida está ordenada}

%---------------------------------------------------------------------
\begin{resultado}{Lema 12.10 · \(\Pairwise\) sobre una concatenación}
\enunciado Para una relación \(R\) y listas \(\ell_1,\ell_2\),
\[
\Pairwise\,R\,(\ell_1\concat\ell_2)\iff
\Pairwise\,R\,\ell_1\ \wedge\ \Pairwise\,R\,\ell_2\ \wedge\
\big(\forall a\in\ell_1,\ \forall b\in\ell_2,\ R\,a\,b\big),
\]
y además \(\Pairwise\,R\,(a::\ell)\iff(\forall b\in\ell,\ R\,a\,b)\wedge\Pairwise\,R\,\ell\).

\palabras Una concatenación está «ordenada por pares» si y solo si cada
trozo lo está y todo elemento del primer trozo está relacionado con todo
elemento del segundo.

\sirve Reduce \(\Sorted(A\concat(p::B))\) a tres comprobaciones sencillas.

\prueba
\textbf{Parte 1 (\(::\)).} Es la definición (P) leída en los dos sentidos.
Si \(\Pairwise\,R\,(a::\ell)\) se obtuvo, fue por la regla de la lista no
vacía, cuyas premisas son las dos condiciones. La regla de la lista vacía no
aplica porque \(a::\ell\neq[\,]\).

\textbf{Parte 2 (\(\concat\)), por inducción sobre \(\ell_1\).} Se escribe
\(C(s,t):=\forall x\in s,\ \forall y\in t,\ R\,x\,y\) y
\(H(s):=\forall z\in s,\ R\,a\,z\).
\begin{itemize}
  \item \(\ell_1=[\,]\): \([\,]\concat\ell_2=\ell_2\), \(\Pairwise\,R\,[\,]\)
    siempre vale y \(C([\,],\ell_2)\) vale por vacuidad.
  \item \(\ell_1=a::r\). Primero dos equivalencias auxiliares (\(\ast\)),
    que solo usan la pertenencia en \(\concat\) y en \(::\):
    \[
    H(r\concat\ell_2)\iff H(r)\wedge H(\ell_2),\qquad
    C(a::r,\ell_2)\iff H(\ell_2)\wedge C(r,\ell_2).
    \]
    Luego la cadena:
    \[
    \begin{aligned}
    \Pairwise\,R\,((a::r)\concat\ell_2)
    &\iff \Pairwise\,R\,(a::(r\concat\ell_2)) &&(1)\\
    &\iff H(r\concat\ell_2)\wedge\Pairwise\,R\,(r\concat\ell_2) &&(2)\\
    &\iff H(r)\wedge H(\ell_2)\wedge\Pairwise\,R\,r &&(3)\\
    &\qquad\quad\wedge\Pairwise\,R\,\ell_2\wedge C(r,\ell_2)\\
    &\iff \big(H(r)\wedge\Pairwise\,R\,r\big)\wedge\Pairwise\,R\,\ell_2 &&(4)\\
    &\qquad\quad\wedge\big(H(\ell_2)\wedge C(r,\ell_2)\big)\\
    &\iff \Pairwise\,R\,(a::r)\wedge\Pairwise\,R\,\ell_2\wedge C(a::r,\ell_2) &&(5)
    \end{aligned}
    \]
    (1) \((a::r)\concat\ell_2=a::(r\concat\ell_2)\); (2) Parte 1;
    (3) la primera equivalencia (\(\ast\)) y la hipótesis inductiva;
    (4) reagrupar conjunciones; (5) Parte 1 y la segunda equivalencia
    (\(\ast\)).
\end{itemize}
\end{resultado}

%---------------------------------------------------------------------
\begin{resultado}{Teorema 12.11 · La salida está ordenada}
\enunciado Para toda lista \(\ell\), \(\Sorted(\qs(\ell))\).

\idea Las dos mitades están ordenadas por inducción. Todo lo de la izquierda
es \(\leq p\) y todo lo de la derecha es \(\geq p\); por transitividad, todo
lo de la izquierda es \(\leq\) todo lo de la derecha.

\prueba Por inducción bien fundada sobre \(\len\ell\). El caso \([\,]\) es
inmediato. Para \(\ell=p::r\), sean \(A:=\qs(L)\) y \(B:=\qs(R)\), de modo que
\(\qs(p::r)=A\concat(p::B)\).
\begin{enumerate}
  \item \textbf{Mitades ordenadas.} Por la hipótesis inductiva (Técnica 2),
    \(\Sorted(A)\) y \(\Sorted(B)\).
  \item \textbf{Cotas.} Por el Corolario 12.9, \(a\in A\iff a\in L\), y
    \(a\in L\) implica \(a\leq p\). Del mismo modo, \(b\in B\) implica
    \(p<b\), luego \(p\leq b\):
    \[
    (\forall a\in A,\ a\leq p)\ \wedge\ (\forall b\in B,\ p\leq b).
    \]
  \item \textbf{El bloque derecho.} Por la Parte 1 del Lema 12.10, con la
    cota derecha y \(\Sorted(B)\), se tiene \(\Sorted(p::B)\).
  \item \textbf{El cruce.} Sean \(a\in A\) y \(c\in p::B\). Si \(c=p\),
    entonces \(a\leq p=c\). Si \(c\in B\), entonces \(a\leq p\leq c\), y por
    transitividad \(a\leq c\).
  \item \textbf{Cierre.} El Lema 12.10, con 1, 3 y 4, da
    \(\Sorted(A\concat(p::B))=\Sorted(\qs(p::r))\).
\end{enumerate}
\end{resultado}

\begin{center}
\begin{tikzpicture}[font=\small]
  \node[draw, rounded corners, fill=green!12, minimum width=3.6cm, minimum height=8mm] (A) {$A=\qs(L)$ (ordenada)};
  \node[draw, rounded corners, fill=blue!10, minimum width=1cm, minimum height=8mm, right=2mm of A] (p) {$p$};
  \node[draw, rounded corners, fill=red!10, minimum width=3.6cm, minimum height=8mm, right=2mm of p] (B) {$B=\qs(R)$ (ordenada)};
  \draw[{Stealth}-{Stealth}, thick] ([yshift=6mm]A.north) -- node[above]{\footnotesize todo $a\in A$ $\leq p$} ([yshift=6mm]p.north);
  \draw[{Stealth}-{Stealth}, thick] ([yshift=6mm]p.north) -- node[above]{\footnotesize $p\leq$ todo $b\in B$} ([yshift=6mm]B.north);
  \draw[{Stealth}-{Stealth}, thick, dashed] ([yshift=13mm]A.north) to[bend left=15] node[above]{\footnotesize todo $a\in A$ $\leq$ todo $b\in B$ (transitividad)} ([yshift=13mm]B.north);
\end{tikzpicture}
\end{center}

%---------------------------------------------------------------------
\begin{resultado}{Teorema 12.12 · Corrección de Quicksort}
\enunciado Para toda lista \(\ell\), \(\qs(\ell)\Perm\ell\) y
\(\Sorted(\qs(\ell))\).

\prueba Es la conjunción de los Teoremas 12.8 y 12.11.
\end{resultado}

%=====================================================================
\section{El ejemplo, comprobado con los lemas}
%=====================================================================

En la raíz del árbol de \([3,1,2,5,4,9]\): \(p=3\), \(r=[1,2,5,4,9]\),
\(L=[1,2]\) y \(R=[5,4,9]\).

\begin{center}
\renewcommand{\arraystretch}{1.3}
{\small
\begin{tabularx}{\linewidth}{@{}p{0.3\linewidth}X@{}}
\toprule
\textbf{Lo que afirma la prueba} & \textbf{En el ejemplo}\\
\midrule
Lema 12.5: \(\len L,\len R\leq\len r\) & \(2,3\leq5<6\)\\
Lema 12.7: \(L\concat R\Perm r\) & \([1,2,5,4,9]\Perm[1,2,5,4,9]\) (aquí incluso son iguales)\\
Hipótesis inductiva & \(A=\qs(L)=[1,2]\), \(B=\qs(R)=[4,5,9]\), ambas ordenadas\\
Cotas & \(1,2\leq3\) y \(3<4,5,9\)\\
Cruce & todo \(\{1,2\}\leq\) todo \(\{3,4,5,9\}\)\\
Lema 12.10 & \([1,2]\concat(3::[4,5,9])=[1,2,3,4,5,9]\) está ordenada\\
Permutación (Teorema 12.8) & cada número aparece una vez en la entrada y una vez en la salida\\
\bottomrule
\end{tabularx}}
\end{center}

\begin{nota}[El Lema 12.7 no siempre es una igualdad]
En la raíz, \(L\concat R\) coincidió con \(r\) por casualidad. En
\(\qs([5,4,9])\) con pivote \(5\): \(r=[4,9]\), \(L=[4]\) y \(R=[9]\), y de
nuevo coinciden. Pero con \(r=[9,4]\) y pivote \(5\) se obtiene
\(L\concat R=[4]\concat[9]=[4,9]\neq[9,4]\), que es solo una permutación.
\end{nota}

%=====================================================================
\section{Qué propiedades del orden se usan, y dónde}
%=====================================================================

\begin{center}
\renewcommand{\arraystretch}{1.3}
{\small
\begin{tabularx}{\linewidth}{@{}p{0.27\linewidth}X@{}}
\toprule
\textbf{Propiedad} & \textbf{Dónde interviene}\\
\midrule
Totalidad & \(\neg(x\leq p)\Rightarrow p<x\): lo que no va a la izquierda es mayor que el pivote. Da la cota derecha del Teorema 12.11.\\
Transitividad & El cruce \(a\leq p\leq c\Rightarrow a\leq c\) del Teorema 12.11.\\
\(p<b\Rightarrow p\leq b\) & Para usar la cota derecha con la relación \(\leq\) de \(\Sorted\).\\
Decidibilidad de \(\leq\) & Para \emph{calcular} los filtros: \(\decidir(x\leq p)\).\\
Igualdad decidible & Para \emph{contar} apariciones, que es como se define \(\Perm\).\\
\bottomrule
\end{tabularx}}
\end{center}

\begin{nota}[Observación propia, no de la tesis]
La antisimetría no interviene en el argumento: el razonamiento funciona igual
para un \emph{preorden total} (por ejemplo, ordenar personas por edad, donde
dos personas distintas pueden tener la misma edad). La formalización trabaja
con orden lineal, que es la hipótesis de la tesis.
\end{nota}

%=====================================================================
\section{Auditoría y alcance}
%=====================================================================

\textbf{Auditoría (§12.10).} La herramienta que lista los axiomas lógicos de
los que depende un resultado reporta, para el teorema de corrección, solo
\emph{extensionalidad proposicional} y \emph{solidez de los cocientes}, sin
ninguna demostración pendiente. A diferencia del Capítulo 11, no aparece el
axioma de elección: la prueba formalizada de la corrección no usa
razonamiento clásico. La definición del algoritmo depende solo de la extensionalidad
proposicional, y \(\qs\) sigue siendo ejecutable: las pruebas se borran al
compilar.

\textbf{Lo que el capítulo no cubre (§12.1).}
\begin{itemize}
  \item Versiones \emph{in situ}, que reordenan dentro del mismo arreglo
    (partición de Hoare), ni el uso de memoria.
  \item La complejidad temporal: \(\Theta(n\log n)\) en promedio y
    \(\Theta(n^2)\) en el peor caso. Se demuestra \emph{qué} calcula el
    algoritmo, no \emph{cuánto} tarda.
\end{itemize}

%=====================================================================
\section{Preguntas probables del tribunal}
%=====================================================================

\begin{description}[style=nextline, leftmargin=0pt, itemsep=6pt]
\item[¿Qué demuestra exactamente el capítulo?]
Que para todo tipo linealmente ordenado y toda lista, Quicksort termina y
devuelve una permutación ordenada de la entrada.

\item[¿Por qué la especificación tiene dos partes?]
Porque cada una sola admite funciones absurdas: \(f(\ell)=[\,]\) siempre
ordena, y \(f(\ell)=\ell\) siempre permuta.

\item[¿Por qué la permutación se define contando y no por pertenencia?]
Porque con pertenencia \([1,1]\) y \([1]\) serían «iguales». Una función que
borrara duplicados pasaría el test. El Corolario 12.9 (pertenencia) es más
débil que el Teorema 12.8.

\item[¿Por qué inducción sobre la longitud y no sobre la lista?]
Porque las llamadas recursivas son sobre \(L\) y \(R\), no sobre la cola. La
inducción estructural solo da hipótesis para la cola.

\item[¿Por qué hay que demostrar la terminación?]
Porque una definición recursiva que no termina puede ser inconsistente
(\(g(n)=g(n)+1\) implicaría \(0=1\)). La prueba garantiza que \(\qs\) es una
función total bien definida.

\item[¿Qué pasa con los elementos repetidos?]
Las copias del pivote van a la izquierda (\(\leq p\)) y se conservan. El
ejemplo \([2,1,2,1]\mapsto[1,1,2,2]\) lo muestra; el Teorema 12.8 lo
garantiza en general.

\item[¿Por qué \(\leq p\) a un lado y el complemento al otro?]
Porque así cada elemento va a exactamente una mitad (Lema 12.7). Si se
usaran \(\leq p\) y \(\geq p\), las copias de \(p\) se duplicarían.

\item[¿Por qué \(\Pairwise\) y no «cada elemento \(\leq\) el siguiente»?]
Son equivalentes por transitividad (Thompson usa la segunda). \(\Pairwise\)
se descompone limpiamente sobre la concatenación (Lema 12.10), que es justo
lo que necesita el paso \(A\concat(p::B)\).

\item[¿Dónde se usa la transitividad?]
En el cruce: todo \(a\in A\) es \(\leq\) todo \(c\in p::B\) porque
\(a\leq p\leq c\).

\item[¿Para qué sirve el Corolario 12.9 si ya está el Teorema 12.8?]
Para obtener las cotas: saber que los elementos de \(\qs(L)\) son los de
\(L\), y por tanto \(\leq p\). Es la forma en que la permutación alimenta el
orden.

\item[¿Hay otra prueba de la permutación?]
Sí, la del Apéndice B.2: contar apariciones y usar
\(c_x(L)+c_x(R)=c_x(r)\).

\item[¿Se demuestra la eficiencia?]
No. El alcance es la corrección funcional, no la complejidad ni la memoria.

\item[¿Usa lógica clásica?]
La prueba formalizada no la usa: la auditoría no reporta el axioma de
elección, del que depende el razonamiento clásico en el sistema. Es una
diferencia con la completitud del Capítulo 11.
\end{description}

%=====================================================================
\appendix
\section{Hoja de referencia}
%=====================================================================

\begin{clave}
\textbf{La prueba completa en tres líneas.} En el caso \(p::r\), con
\(L\) los \(\leq p\) y \(R\) los \(>p\) de \(r\):
\begin{enumerate}
  \item \(\len L,\len R<\len{p::r}\) (Lema 12.5), así que la hipótesis
    inductiva vale para \(L\) y \(R\).
  \item \textbf{Permutación:}
    \(\qs(L)\concat(p::\qs(R))\Perm L\concat(p::R)\Perm p::(L\concat R)\Perm p::r\).
  \item \textbf{Orden:} \(A,B\) ordenadas; \(A\leq p\leq B\); el Lema 12.10
    pega las tres piezas.
\end{enumerate}
\end{clave}

\textbf{Definiciones clave.}
\[
\qs(p::r)=\qs(L)\concat(p::\qs(R)),
\]
\[
\Pairwise\,R\,(a::r)\iff(\forall x\in r,\ R\,a\,x)\wedge\Pairwise\,R\,r,
\]
\[
\ell_1\Perm\ell_2\iff\forall a,\ \cuenta(a,\ell_1)=\cuenta(a,\ell_2).
\]

\textbf{Los tres datos que más se preguntan:}
\begin{itemize}
  \item La inducción es sobre la \emph{longitud}, porque \(L\) y \(R\) no son
    la cola.
  \item La permutación cuenta multiplicidades; la pertenencia sola no basta.
  \item El orden de la salida sale de las cotas respecto del pivote más la
    transitividad.
\end{itemize}

\end{document}
